\section{Conclusion}

The study of Chapter 8 shows that malicious software has changed from a small technical problem into a professional criminal tool. As malware becomes smarter, our ways of defending must also improve.

\subsection{Final Summary of Key Lessons}

Through our research, we can conclude three main things:

\begin{itemize}
  \item Technology is not the only answer: We can have the best firewalls, but Social Engineering (tricking people) is still the easiest way for hackers to get in. Educating users (User Awareness) is the most important layer of defense [1].
  \item Protecting the CIA Triad: Every attack we studied targets the pillars of security. Spyware breaks Confidentiality, Viruses break Integrity, and Ransomware destroys Availability. A good security plan must protect all three [1].
  \item Focus on Resilience: Since threats like "Fileless" and "AI malware" are hard to stop, organizations must focus on Cyber Resilience. This means being ready to recover quickly using a strong 3-2-1 Backup Strategy and a clear Incident Response plan.
\end{itemize}

\subsection{The Reality in Vietnam}

The 2024-2025 Ransomware attacks in Vietnam, which hit companies like VNDIRECT and PVOIL, prove that the theories in Chapter 8 are very real. These attacks caused huge losses and showed that Vietnamese businesses must implement Defense in Depth immediately to survive in the digital age.

\subsection{Final Words}

Malware will never stop evolving. However, by staying proactive with Patch Management, using modern AI-based defense tools, and constantly teaching users to be careful, we can stay safe. Security is not a one-time job; it is a continuous journey of learning and adapting.
